\documentclass{sbc2023}

\usepackage{graphicx}
\usepackage{fontspec}
\usepackage{color}
\usepackage{hyperref}
\usepackage{url}
\usepackage{aas_macros}
\usepackage[bottom]{footmisc}
\usepackage{amsmath}
\usepackage{amssymb}
\usepackage{tabularx}
\usepackage{booktabs}
\usepackage[portuguese,linesnumbered,ruled,vlined]{algorithm2e}
\usepackage{multirow}
\usepackage{multicol}

\setcitestyle{square}

\definecolor{orcidlogo}{rgb}{0.37,0.48,0.13}
\definecolor{darkblue}{rgb}{0.0,0.0,0.0}
\hypersetup{colorlinks,breaklinks,
            linkcolor=darkblue,urlcolor=darkblue,
            anchorcolor=darkblue,citecolor=darkblue}

\jid{JBCS}
\jtitle{Journal of the Brazilian Computer Society, 202X, XX:1, }
\doi{10.5753/jbcs.202X.XXXXXX}
\copyrightstatement{This work is licensed under a Creative Commons Attribution 4.0 International License}
\jyear{202X}

\title[Sistemas de Recomenda{\c c}{\~a}o e Sa{\'u}de Mental]{SISTEMAS DE RECOMENDA{\c C}{\~A}O E SA{\'U}DE MENTAL NO CAPITALISMO DE VIGIL{\^A}NCIA: A CURADORIA INTENCIONAL COMO INSTRUMENTO DE MITIGA{\c C}{\~A}O DA SOBRECARGA COGNITIVA}

\author[Melo, Farias \& Marciano]{
\affil{\textbf{Fellype Samuel}~\href{https://orcid.org/0000-0000-0000-0000}{\textcolor{orcidlogo}{\aiOrcid}}~[~\textbf{FAETERJ-RIO de Quintino, RJ, Brasil}~]}
\affil{\textbf{Felipe Farias}~\href{https://orcid.org/0000-0000-0000-0000}{\textcolor{orcidlogo}{\aiOrcid}}~[~\textbf{FAETERJ-RIO de Quintino, RJ, Brasil}~]}
\affil{\textbf{Ricardo Marciano}~\href{https://orcid.org/0000-0000-0000-0000}{\textcolor{orcidlogo}{\aiOrcid}}~[~\textbf{FAETERJ-RIO de Quintino, RJ, Brasil}~]}
}

\SetKwInput{KwData}{Entrada}
\SetKwInput{KwResult}{Sa{\'\i}da}
\SetKw{KwRetorna}{retornar}
\SetKw{KwSe}{se}
\SetKw{KwSenao}{sen{\~a}o}
\SetKw{KwParaCada}{para cada}
\SetKw{KwEnquanto}{enquanto}

\begin{document}

\begin{frontmatter}
\maketitle

\begin{mail}
fellype.24104708360042@faeterj-rio.edu.br, felipe.23204708360023@faeterj-rio.edu.br, ricardo.marciano@faeterj-rio.edu.br
\end{mail}

\begin{dates}
\small{\textbf{Received:} DD Month YYYY~~~$\bullet$~~~\textbf{Accepted:} DD Month YYYY~~~$\bullet$~~~\textbf{Published:} DD Month YYYY}
\end{dates}

\begin{abstract}
\textbf{Resumo.~}
\noindent \textbf{Introdução:} A economia da atenção e o capitalismo de vigilância impulsionam sistemas de recomendação focados na maximização do engajamento de curto prazo, associados à fadiga atencional e a impactos negativos na saúde mental. \textbf{Objetivo:} Propor e validar uma arquitetura algorítmica alternativa que priorize a atenção sustentável e a preservação da utilidade cognitiva. \textbf{Metodologia:} A arquitetura substitui a recompensa por engajamento bruto por um Score de Qualidade Multiobjetivo, modela a interação via Processos de Hawkes decompostos (Sistema 1 vs. Sistema 2) e introduz uma EDO de reserva cognitiva avaliada localmente por Edge AI. A validação foi realizada por Simulação Baseada em Agentes (ABM) com N = 1.000 agentes. \textbf{Resultados:} A arquitetura proposta preserva a reserva cognitiva ($d = 3.25; p < 10^{-300}$) e reduz a divergência entrópica do usuário, demonstrando robustez em análises de sensibilidade. \textbf{Conclusão:} Os resultados oferecem um caminho técnico pragmático para sistemas de recomendação eticamente alinhados, provando que é possível otimizar a utilidade sustentável sem prescindir da eficiência técnica.
\end{abstract}

\begin{keywords}
Sistemas de Recomenda{\c c}{\~a}o, Sa{\'u}de Mental, Capitalismo de Vigil{\^a}ncia, Processos de Hawkes, Intelig{\^e}ncia Humana.
\end{keywords}

\end{frontmatter}

\section{Introdução}
Nas últimas duas décadas, as redes sociais tornaram-se o centro da vida social, da informação e do entretenimento humano. Contudo, avolumam-se as discussões empíricas sobre como tais ferramentas afetam a saúde mental, a capacidade de concentração e a qualidade do consumo online. A crise de atenção contemporânea não se resume a uma falha de disciplina individual, mas deriva de estruturas econômicas e tecnológicas que induzem um déficit cognitivo sistêmico \citep{hari2022}. As plataformas são projetadas para maximizar a retenção por meio de ciclos de recompensa intermitente, instigando a verificação compulsiva.

Esse cenário opera sob a égide do ``capitalismo de vigilância'' \citep{zuboff2019}, modelo em que cada interação é quantificada para alimentar algoritmos que mapeiam vulnerabilidades. \citet{citton2017} descreve uma ``ecologia da atenção'' na qual algoritmos bombardeiam o indivíduo com estímulos rápidos. A disputa converte a atenção em um ativo econômico. Os efeitos recaem de forma acentuada sobre jovens, com associações entre o uso intensivo de algoritmos e sintomas emocionais adversos \citep{costello2023}. Relatórios institucionais corroboram que sistemas não supervisionados estão atrelados à exaustão e ansiedade digitais \citep{hai2021}.

A literatura sociotécnica descreve as plataformas digitais como infraestruturas centrais de um regime de capitalismo de vigilância, no qual dados comportamentais são extraídos para a modulação de condutas e predição comercial \citep{zuboff2019}. Especialistas apontam que a otimização calcada puramente no engajamento subverte os interesses do indivíduo \citep{harris2019}. \citet{citton2017} documenta uma ecologia em que a escassez de foco é produzida de modo algorítmico e artificial por interações massivas de reatividade, enquanto métricas corroboram como essa exaustão consome componentes significativos do processamento deliberativo dos agentes no médio prazo \citep{allcott2021}.

De forma a mitigar tais externalidades, consolida-se o campo de \textit{Value-Aligned RecSys}, em que pesquisadores modelam algoritmos não baseados em simples predições brutas, contendo arranjos paralelos para o engajamento contínuo, a justiça ou a pluralidade em distribuição de ganhos na rede de contato do usuário \citep{stray2021}. Por meio do \textit{reinforcement learning} aplicado ao ambiente sociopolítico da plataforma, avalia-se que recompensas orientadas estritamente a fluxos ininterruptos fragmentam as teias colaborativas dos usuários \citep{carvalho2023}.

Correlato a isso, o bem-estar digital costuma ser abordado pela adoção de dispositivos externos de autogestão, como dashboards e aplicativos de bloqueio \citep{forest, rescuetime}, cuja efetividade é limitada por dependerem exclusivamente do autocontrole \citep{lyngs2019}. Soluções efetivas advêm de mecanismos coercitivos projetados via fricções intencionais \citep{kushlev2022}.

Os elos da economia comportamental com o design de sistemas fornecem métricas robustas. A segregação dual entre Sistema 1 (automático, impulsivo) e Sistema 2 (reflexivo, deliberativo) \citep{kahneman2011} e os limiares de sobrecarga cognitiva \citep{paas2023} têm direcionado novos desenvolvimentos. Interseccionando essa restrição fisiológica com Processos de Hawkes, é possível discernir se a reiteração da utilização baseia-se em excitação residual curta ou em utilidade orgânica prolongada \citep{rizoiu2017, agarwal2024}.

A convergência dessas discussões evidencia uma lacuna crítica: enquanto a literatura mapeia os riscos do capitalismo de vigilância e propõe modelos de alinhamento de valores, a implementação prática de ferramentas que integrem a regulação cognitiva em tempo real ainda é incipiente. O projeto \textit{Be-Productive} preenche essa lacuna ao transpor a teoria da ecologia da atenção \citep{citton2017} e os modelos de Processos de Hawkes \citep{rizoiu2017} para uma arquitetura funcional de \textit{Edge AI}, automatizando a preservação da reserva cognitiva através de travamentos algorítmicos baseados na dinâmica de sistemas.

Diante desta conjuntura, formula-se a seguinte questão de pesquisa: \textit{É viável estruturar o núcleo de sistemas de recomendação para otimizar a utilidade sustentável e a intenção declarada do usuário, mitigando a dependência algorítmica?}

Este artigo propõe uma arquitetura algorítmica voltada à atenção sustentável. A Seção 2 descreve a metodologia e a arquitetura proposta; a Seção 3 apresenta os resultados e a discussão; e a Seção 4 consolida as conclusões.

\section{Metodologia}
\subsection{A Matemática do Engajamento Imediato}
Nos sistemas comerciais, a plataforma atua primariamente na maximização da interação de curto prazo. Seja $\mathcal{U}$ o conjunto de usuários e $\mathcal{I}$ o de itens, a política $\pi$ otimiza:
\begin{equation} \label{eq:engajamento}
\max_{\pi} \mathbb{E}_{\pi} \left[ \sum_{t=0}^{T} \gamma^t R_E(u, i, t) \right]
\end{equation}
Onde $\gamma \in (0, 1]$ é o fator de desconto temporal e $R_E$ representa a recompensa imediata (ex.: clicks). Esse ranqueamento está demonstrado no Algoritmo 1.

\begin{algorithm}[h]
\caption{Lógica Tradicional de Ranqueamento de Feed}
\KwData{$post \in \mathcal{I}$, $usuario \in \mathcal{U}$}
\KwResult{Pontuação baseada em interação iminente}
$pontuacao \gets (post.curtidas \times w_{curtida}) + (post.compartilhamentos \times w_{compartilhamento})$\;
\If{$pontuacao > LIMIAR\_VIRALIDADE$}{
    $injetar\_no\_feed\_global(post)$\;
}
\KwRetorna{$pontuacao$}\;
\end{algorithm}

\subsection{Dicotomia Cognitiva e Processos de Hawkes}
Para superar as mazelas inerentes e estruturais observadas nos sistemas de recomendação convencionais, propomos uma modelagem matemática da interação usuário-plataforma baseada em Processos de Hawkes, que permite discernir quantitativamente dois regimes comportamentais fundamentais: o consumo impulsivo (reativo) e o consumo intencional (deliberativo). Esta distinção operacionaliza, no domínio algorítmico, a teoria dos sistemas duais de \citet{kahneman2011}, separando o Sistema 1 (automático, rápido, emocional) do Sistema 2 (reflexivo, lento, racional).

\subsubsection{Fundamentos dos Processos de Hawkes}
Um Processo de Hawkes é um processo pontual autorregressivo no qual a ocorrência de eventos passados aumenta a probabilidade de eventos futuros, modelando assim excitação temporal e dependência histórica. Formalmente, a intensidade condicional $\lambda(t)$ (taxa instantânea de ocorrência de eventos no instante $t$) é definida como:
\begin{equation}
\lambda(t) = \mu + \sum_{t_i<t} \phi(t - t_i)
\end{equation}
onde $\mu > 0$ é a taxa base e $\phi(\cdot)$ é o kernel de excitação, que descreve como cada evento passado contribui para a intensidade atual. Tipicamente, adota-se um kernel exponencial decrescente $\phi(\Delta) = \alpha e^{-\beta \Delta}$.

\subsubsection{Decomposição em Sistemas 1 e 2}
A equação central de intensidade é:
\begin{equation} \label{eq:hawkes}
\lambda(t) = \mu + \sum_{t_k \in \mathcal{K}, t_k < t} \alpha_1 e^{-\beta_1(t-t_k)} + \sum_{t_m \in \mathcal{M}, t_m < t} \alpha_2 e^{-\beta_2(t-t_m)}
\end{equation}
onde $\mathcal{K}$ é o conjunto de eventos do Sistema 1 (impulsivos) e $\mathcal{M}$ o do Sistema 2 (intencionais). Parâmetros $\alpha_1, \beta_1$ representam alta excitação e rápido decaimento, enquanto $\alpha_2, \beta_2$ refletem persistência e profundidade.

\subsection{Score de Qualidade e Filtragem Multiobjetivo}
A política proposta substitui o engajamento imediato por um Score de Qualidade $Q(i, u)$ que integra relevância personalizada, feedback explícito e segurança de conteúdo.
\begin{equation} \label{eq:score}
Q(i, u) = \left[ \sum_{k=1}^{K} w_k \cdot f_k(i, u) + \sum_{j=1}^{J} \theta_j \cdot E_j(i, u) \right] \times \min_m (1 - P_m(i))
\end{equation}
O fator $\min_m (1 - P_m(i))$ atua como uma penalidade mínima contra nocividade (clickbait, discurso de ódio, estímulos compulsivos).

\subsection{EDO de Reserva e Edge AI}
A variável $R(t)$ representa a Reserva Cognitiva, um recurso limitado que é drenado por interações impulsivas e restaurado homeostaticamente. A dinâmica é regida por:
\begin{equation} \label{eq:edo}
\frac{dR(t)}{dt} = S(R, t) - L(X, t)
\end{equation}
onde $L(X, t) = \kappa_1 \cdot v_{scroll}(t) + \kappa_2 \cdot \nu_{alt}(t)$ quantifica o esgotamento por sensores comportamentais processados via \textit{Edge AI}. Quando $R(t)$ atinge um limiar crítico, o sistema dispara fricções positivas para induzir a recuperação.

\section{Resultados e Discussão}
A validação via ABM com $N = 1.000$ agentes demonstrou eficácia profunda:
\begin{itemize}
    \item \textbf{Significância Estatística ($p < 10^{-300}$):} A melhoria na preservação da reserva é robusta e não aleatória.
    \item \textbf{Tamanho de Efeito (Cohen's $d = 3.25$):} Magnitude radical de transformação na experiência atencional.
    \item \textbf{Divergência KL:} Redução drástica na distância entre intenção e alocação real de tempo, aproximando-se de zero no modelo proposto.
\end{itemize}

A arquitetura atua como um ``airbag cognitivo'', automatizando a preservação da agência humana. A execução \textit{on-device} (Edge AI) garante a privacidade \textit{by design}, alinhando a técnica aos preceitos éticos contemporâneos.

\section{Conclusão}
O projeto \textit{Be-Productive} estabelece um novo paradigma para a Interação Humano-Computador. Comprovamos que é factível e tecnicamente eficiente abdicar da predação atencional em favor da agência e saúde mental do usuário. O alinhamento de valores não é apenas um imperativo filosófico, mas uma factibilidade de engenharia.

\bibliographystyle{apalike}
\bibliography{referencias}

\end{document}
